\documentclass[11pt]{article}
\newcommand{\DocRunningHead}{Intermediate Algebra}
\newcommand{\DocShortTitle}{Chapters 1--2 Quiz --- Solutions}
% Shared preamble for the Chapter 1-2 quiz and its solutions.
\usepackage[letterpaper,margin=1in,top=0.9in,bottom=0.9in]{geometry}
\usepackage{amsmath,amssymb}
\usepackage[T1]{fontenc}
\usepackage{newpxtext,newpxmath}
\usepackage{enumitem}
\usepackage{fancyhdr}
\usepackage{xcolor}
\usepackage{titlesec}
\usepackage[hidelinks]{hyperref}

\definecolor{rule}{gray}{0.45}
\definecolor{faint}{gray}{0.35}

% --- running head / foot ---
\pagestyle{fancy}
\fancyhf{}
\renewcommand{\headrulewidth}{0.4pt}
\renewcommand{\footrulewidth}{0pt}
\fancyhead[L]{\footnotesize\color{faint}\DocRunningHead}
\fancyhead[R]{\footnotesize\color{faint}\DocShortTitle}
\fancyfoot[C]{\footnotesize\color{faint}Page \thepage\ of \pageref{LastPage}}
\usepackage{lastpage}

% --- part headings ---
\titleformat{\section}{\large\bfseries}{}{0pt}{}
\titlespacing*{\section}{0pt}{16pt}{6pt}
\newcommand{\parthead}[2]{%
  \section*{\rule[-2pt]{3pt}{14pt}\hspace{6pt}#1\hfill\normalsize\mdseries\color{faint}#2}%
  \vspace{-4pt}\noindent\textcolor{rule}{\rule{\linewidth}{0.4pt}}\vspace{2pt}}

% --- title block ---
\newcommand{\quiztitle}[3]{%
  \begin{center}
    {\footnotesize\color{faint}\MakeUppercase{#1}}\\[3pt]
    {\LARGE\bfseries #2}\\[4pt]
    {\small\color{faint}#3}
  \end{center}
  \vspace{2pt}\noindent\textcolor{rule}{\rule{\linewidth}{1pt}}\vspace{6pt}}

% --- name / date / score rules ---
\newcommand{\namedate}{%
  \noindent\small
  Name:\ \rule[-4pt]{0.36\linewidth}{0.4pt}\hfill
  Date:\ \rule[-4pt]{0.16\linewidth}{0.4pt}\hfill
  Score:\ \rule[-4pt]{0.09\linewidth}{0.4pt}\,/\,100\par\vspace{10pt}}

% --- question list ---
\newlist{questions}{enumerate}{1}
\setlist[questions]{label=\textbf{\arabic*.},leftmargin=*,labelsep=8pt,itemsep=0pt,parsep=4pt}

% Vertical filler between questions. Using \vfill (rather than a fixed \vspace)
% means work space is distributed to fill each page and is never silently
% discarded at a page break.
\newcommand{\work}{\vfill}

% solution formatting
\newcommand{\answer}[1]{\fbox{$\displaystyle #1$}}
\newcommand{\answertext}[1]{\fbox{\parbox{0.92\linewidth}{#1}}}
\newenvironment{solution}
  {\par\vspace{2pt}\begingroup\small\setlength{\leftskip}{12pt}%
   \noindent\textcolor{rule}{\rule[-1pt]{1.2pt}{\baselineskip}}\hspace{4pt}\ignorespaces}
  {\par\endgroup\vspace{8pt}}
\newcommand{\source}[1]{{\footnotesize\color{faint}\textit{#1}}}


\begin{document}

\quiztitle{Intermediate Algebra --- Instructor Copy}{Chapters 1 \& 2 Quiz: Solutions}
  {Full worked solutions, with the source problem cited for each question}

\noindent\small
\textbf{Grading note.}\ Each question is worth 10 points. The boxed value is the answer;
the lines above it are the work students are expected to show. Where a problem has a
standard trap --- an extraneous root, a strict versus non-strict domain endpoint, an
excluded value in a range --- it is flagged in italics so partial credit can be applied
consistently.\par

\vspace{4pt}

\parthead{Part I --- Chapter 1}{Basic Techniques for Solving Equations}

\begin{questions}

\item \source{Review Problem 1.23(b), \S1.1}\\
  Solve $x - 1 \le 3x + 2 \le 2x + 6$; write the answer in interval notation.
  \begin{solution}
    A double inequality is two inequalities at once; solve each separately and intersect.

    Left: $x - 1 \le 3x + 2$. Subtracting $x$ from both sides gives $-1 \le 2x + 2$,
    so $-3 \le 2x$, and $x \ge -\tfrac{3}{2}$.

    Right: $3x + 2 \le 2x + 6$. Subtracting $2x$ gives $x + 2 \le 6$, so $x \le 4$.

    Both must hold, so the solution is the overlap.
    \[ \answer{\left[-\tfrac{3}{2},\ 4\right]} \]
    \textit{Note: no multiplication or division by a negative occurs here, so no inequality
    sign flips. Both endpoints are included because both inequalities are non-strict.}
  \end{solution}

\item \source{Review Problem 1.19(a), \S1.3}\\
  Solve the system $4x + 5y = 43$, $9x - 2y = 57$.
  \begin{solution}
    Eliminate $y$: multiply the first equation by 2 and the second by 5 so the
    $y$-coefficients become $+10$ and $-10$.
    \[
      \begin{aligned}
        8x + 10y &= 86,\\
        45x - 10y &= 285.
      \end{aligned}
    \]
    Adding gives $53x = 371$, so $x = 7$. Substituting into $4x + 5y = 43$ gives
    $28 + 5y = 43$, so $5y = 15$ and $y = 3$.

    Check in the second original equation: $9(7) - 2(3) = 63 - 6 = 57$. \checkmark
    \[ \answer{(x,y) = (7,3)} \]
  \end{solution}

\item \source{Review Problem 1.20, \S1.2}\\
  Jeff is 4 times older than his daughter. Five years ago he was 9 times older than his
  daughter. How old is his daughter?
  \begin{solution}
    Let $d$ be the daughter's age now; then Jeff's age is $4d$. Five years ago their ages
    were $d - 5$ and $4d - 5$, and the second was 9 times the first:
    \[ 4d - 5 = 9(d - 5). \]
    Expanding gives $4d - 5 = 9d - 45$, so $40 = 5d$ and $d = 8$.

    Check: the daughter is 8 and Jeff is 32; five years ago they were 3 and 27, and
    $27 = 9 \cdot 3$. \checkmark
    \[ \answer{\text{The daughter is } 8 \text{ years old.}} \]
    \textit{Note: students who set up $4d - 5 = 9d - 5$ (forgetting to age the daughter
    down as well) get $d = 0$; the check is what catches this.}
  \end{solution}

\item \source{Review Problem 1.24, \S1.2}\\
  Find all ordered pairs $(x,y)$ with $\sqrt{x} + \sqrt{y} = 7$ and
  $3\sqrt{x} - 4\sqrt{y} = -14$.
  \begin{solution}
    The system is linear in $\sqrt{x}$ and $\sqrt{y}$, so substitute
    $u = \sqrt{x}$ and $v = \sqrt{y}$, where $u \ge 0$ and $v \ge 0$:
    \[
      \begin{aligned}
        u + v &= 7,\\
        3u - 4v &= -14.
      \end{aligned}
    \]
    From the first, $u = 7 - v$. Substituting into the second gives
    $3(7 - v) - 4v = -14$, so $21 - 7v = -14$, hence $v = 5$ and $u = 2$.

    Both are non-negative, so both are legitimate square roots. Squaring,
    $x = u^2 = 4$ and $y = v^2 = 25$.

    Check: $\sqrt{4} + \sqrt{25} = 2 + 5 = 7$ and $3(2) - 4(5) = -14$. \checkmark
    \[ \answer{(x,y) = (4,25)} \]
    \textit{Note: a negative value of $u$ or $v$ would have to be discarded, since
    $\sqrt{\phantom{x}}$ denotes the non-negative root. Here neither is negative, so
    exactly one ordered pair survives.}
  \end{solution}

\item \source{Review Problem 1.21(a), \S1.4}\\
  Solve $8x + y - z = 46$, \; $3x + 4y - 2z = 27$, \; $4x - y + z = 14$.
  \begin{solution}
    The first and third equations have opposite $y$- and $z$-coefficients, so adding them
    eliminates both variables at once:
    \[ (8x + y - z) + (4x - y + z) = 46 + 14 \quad\Longrightarrow\quad 12x = 60, \]
    so $x = 5$.

    Substituting $x = 5$ into the third equation gives $20 - y + z = 14$, so $y - z = 6$.

    Substituting $x = 5$ into the second gives $15 + 4y - 2z = 27$, so $4y - 2z = 12$,
    that is, $2y - z = 6$.

    Subtracting $y - z = 6$ from $2y - z = 6$ gives $y = 0$, and then $z = y - 6 = -6$.

    Check: $8(5) + 0 - (-6) = 46$ \checkmark, \;
    $3(5) + 0 - 2(-6) = 27$ \checkmark, \;
    $4(5) - 0 + (-6) = 14$ \checkmark.
    \[ \answer{(x,y,z) = (5,\,0,\,-6)} \]
    \textit{Note: spotting that equations 1 and 3 add to eliminate two variables is the
    intended shortcut, but any correct elimination or substitution order earns full credit.}
  \end{solution}

\end{questions}

\parthead{Part II --- Chapter 2}{Functions Review}

\begin{questions}[resume]

\item \source{Review Problem 2.23(d), \S2.1}\\
  Find the domain and range of $S(x) = \dfrac{12x-9}{6-9x}$.
  \begin{solution}
    \textbf{Domain.} The only restriction is that the denominator cannot vanish:
    $6 - 9x = 0$ when $x = \tfrac{2}{3}$.
    \[ \answer{\text{Domain: } \left(-\infty,\tfrac{2}{3}\right) \cup \left(\tfrac{2}{3},\infty\right)} \]

    \textbf{Range.} Use the chapter's strategy: set the function equal to a new variable
    and solve for $x$. Writing $y = \dfrac{12x-9}{6-9x}$ and clearing the denominator,
    \[ y(6-9x) = 12x - 9 \quad\Longrightarrow\quad 6y - 9xy = 12x - 9. \]
    Collect the $x$ terms on one side: $6y + 9 = 12x + 9xy = x(12 + 9y)$, so
    \[ x = \frac{6y+9}{9y+12}. \]
    This produces a valid input $x$ for every $y$ except the one making the denominator
    zero: $9y + 12 = 0$, i.e. $y = -\tfrac{4}{3}$. That value of $y$ is therefore never
    attained, and every other value is.
    \[ \answer{\text{Range: } \left(-\infty,-\tfrac{4}{3}\right) \cup \left(-\tfrac{4}{3},\infty\right)} \]
    \textit{Note: $-\tfrac43$ is the ratio of the leading coefficients, $12/(-9)$ --- the
    horizontal asymptote. Answers of ``all real numbers'' for the range are the common
    error and should lose the range half of the credit.}
  \end{solution}

\item \source{Review Problem 2.24(a), \S2.1}\\
  Find the domain of $f(x) = \dfrac{1}{\sqrt{2x-5}} + \sqrt{9-3x}$.
  \begin{solution}
    Each term imposes its own restriction, and $x$ must satisfy both.

    First term: $\sqrt{2x-5}$ must be defined \emph{and} nonzero, since it sits in a
    denominator. So $2x - 5 > 0$, giving $x > \tfrac{5}{2}$ --- a \emph{strict} inequality.

    Second term: $9 - 3x \ge 0$, giving $x \le 3$.

    Intersecting the two conditions:
    \[ \answer{\left(\tfrac{5}{2},\ 3\right]} \]
    \textit{Note: the endpoint asymmetry is the whole point of the problem. Writing
    $\left[\tfrac52, 3\right]$ (treating the first radical as if it were not in a
    denominator) is the expected error.}
  \end{solution}

\item \source{Review Problem 2.35, \S2.2}\\
  How can the graph of $f(x)$ be used to produce the graph of $y = |f(x)|$?
  \begin{solution}
    By the definition of absolute value,
    \[
      |f(x)| =
      \begin{cases}
        f(x), & \text{if } f(x) \ge 0,\\
        -f(x), & \text{if } f(x) < 0.
      \end{cases}
    \]
    So wherever the graph of $f$ is on or above the $x$-axis, the graph of $|f|$ is
    identical to it; wherever the graph of $f$ is below the $x$-axis, the output changes
    sign, which reflects that piece across the $x$-axis.
    \[ \answertext{Keep every part of the graph of $f$ that lies on or above the
      $x$-axis, and reflect every part that lies below the $x$-axis across the $x$-axis.} \]
    \textit{Full credit requires both halves of the description plus a reason. The
    resulting graph never dips below the $x$-axis, and it meets the original graph exactly
    at the $x$-intercepts of $f$, where the two pieces join.}
  \end{solution}

\item \source{Review Problem 2.25, \S2.3}\\
  With $f(x) = \dfrac{4x}{x+2}$ and $g(x) = \dfrac{2x}{x+4}$, find $f(g(x))$.
  \begin{solution}
    Substitute $g(x)$ into $f$ wherever $x$ appears:
    \[
      f(g(x)) = \frac{4 \cdot \dfrac{2x}{x+4}}{\dfrac{2x}{x+4} + 2}
              = \frac{\dfrac{8x}{x+4}}{\dfrac{2x}{x+4} + 2}.
    \]
    Multiply the numerator and the denominator by $x + 4$ to clear the inner fractions:
    \[
      f(g(x)) = \frac{8x}{2x + 2(x+4)} = \frac{8x}{4x + 8}.
    \]
    Dividing numerator and denominator by 4,
    \[ \answer{f(g(x)) = \frac{2x}{x+2}} \]
    \textit{Note: strictly, this composition is undefined at $x = -4$ (outside the domain
    of $g$) and at $x = -2$ (where $g(x) + 2 = 0$), even though the simplified formula
    hides the first of these. Mentioning either restriction is a bonus, not a requirement.}
  \end{solution}

\item \source{Review Problem 2.37, \S2.4 \;(Source: AHSME)}\\
  $f(x) = \dfrac{cx}{2x+3}$ satisfies $f(f(x)) = x$ for all real $x$ except $-\tfrac{3}{2}$.
  Find $c$.
  \begin{solution}
    Compute the composition, multiplying numerator and denominator by $2x+3$:
    \[
      f(f(x)) = \frac{c \cdot \dfrac{cx}{2x+3}}{2 \cdot \dfrac{cx}{2x+3} + 3}
              = \frac{c^2 x}{2cx + 3(2x+3)}
              = \frac{c^2 x}{(2c+6)x + 9}.
    \]
    Setting this equal to $x$ and clearing the denominator gives
    \[ c^2 x = x\big[(2c+6)x + 9\big] = (2c+6)x^2 + 9x, \]
    so $(2c+6)x^2 + (9 - c^2)x = 0$ must hold for \emph{every} allowed $x$. A polynomial
    that is identically zero has all coefficients zero:
    \[ 2c + 6 = 0 \quad\text{and}\quad 9 - c^2 = 0. \]
    The first gives $c = -3$; the second gives $c = \pm 3$. Only $c = -3$ satisfies both.
    \[ \answer{c = -3} \]
    Check: with $f(x) = \dfrac{-3x}{2x+3}$,
    \[
      f(f(x)) = \frac{-3 \cdot \dfrac{-3x}{2x+3}}{2 \cdot \dfrac{-3x}{2x+3} + 3}
              = \frac{9x}{-6x + 6x + 9} = \frac{9x}{9} = x. \;\checkmark
    \]
    \textit{Note: $f$ being its own inverse is what $f(f(x)) = x$ says. Students who
    substitute a single convenient value of $x$ (say $x = 1$) can find a candidate for $c$
    but have not shown the identity holds for all $x$; that is worth partial credit only
    if they verify the candidate afterwards, as above.}
  \end{solution}

\end{questions}

\vspace{6pt}
\noindent\textcolor{rule}{\rule{\linewidth}{0.4pt}}\\[2pt]
\source{All problems are taken from Richard Rusczyk, David Patrick, and Ravi Boppana,
\textit{Art of Problem Solving: Intermediate Algebra}, Chapters 1--2 (Review Problems).
Every answer above was verified symbolically with SymPy.}

\end{document}
